% Copyright (c) 2025 ArkLib Contributors. All rights reserved.
% Released under Apache 2.0 license as described in the file LICENSE.
% Authors: Poulami Das (Least Authority)

\section{STIR}

\subsection{Tools for Reed Solomon codes}
\subsubsection{Random linear combination as a proximity generator}\label{sec:proximity_gap}

\begin{theorem}\label{thm:proximity_gap}
\lean{STIR.proximity_gap}
\leanok
\uses{def:distance_from_code,def:reed_solomon_code,def:list_close_codewords}
    Let $\code := \rscode[\field,\evaldomain,\degree]$ be a Reed Solomon code with $\degree\ge1$ and rate $\rate:=\frac{\degree}{|\evaldomain|}$ and let $B^*(\rate):=\sqrt{\rate}$. For every $\distance \in (0,1 - B^*(\rate))$ and functions $f_0,\ldots,f_{m-1} : \evaldomain \to \field$, if
    \[
    \Pr_{r\leftarrow\field}\!\Bigl[
      \Delta\Bigl(\sum_{j=0}^{m-1} r^j\cdot {f_j},\rscode[\field,\evaldomain,\degree]\Bigr)
      \le \delta
    \Bigr]>\err^*(\degree,\rate,\distance,m),
    \]
    then there exists a subset $S \subseteq \evaldomain$ with $|S| \ge (1 - \delta)\cdot|\evaldomain|$,
    and for every $i \in \{0,\ldots,m-1\}$, there exists $u \in \rscode[\field,\evaldomain,\degree]$ such that $f_i(S) = u(S)$.

    \medskip

    \noindent
    Above, $\err^*(\degree,\rate,\distance,m)$ is defined as follows:
    \begin{itemize}
        \item if $\distance \in \left(0,\frac{1-\rate}{2}\right]$ then
            \[
                \err^*(\degree,\rate,\distance,m)=\frac{(m-1)\cdot \degree}{\rate\cdot|\field|}
            \]
        \item if $\distance \in \Bigl(\frac{1-\rate}{2}, 1-\sqrt{\rate}\Bigr)$ then
        \[
            \err^*(\degree,\rate,\distance,m)=\frac{(m-1)\cdot {\degree}^2}{|\field|\cdot{\Bigl(2\cdot\min\{1-\sqrt{\rate}-\distance,\frac{\sqrt{\rate}}{20}\}\Bigr)}^7}
        \]
    \end{itemize}
\end{theorem}

\begin{proof}
\leanok
For an empty family, take the whole evaluation domain. Otherwise write $m=k+1$ and apply
correlated agreement for affine curves from \cite{BCIKS20}. In the admissible distance range,
the error above equals $k$ times the correlated-agreement error bound, in both the unique-decoding
and Johnson regimes. The resulting joint agreement gives the common set $S$.
\end{proof}

\subsubsection{Univariate Function Quotienting}\label{sec:quotienting}

In the following, we start by defining the \emph{quotient} of a univariate function.
\begin{definition}\label{def:quotient}
\lean{Quotienting.funcQuotient}
\leanok
    Let $f:\evaldomain\to\field$ be a function, $S\subseteq\field$ be a finite set, and $\mathsf{Ans},\mathsf{Fill}: S\rightarrow\field$ be functions. Let $\hat{\mathsf{Ans}}\in\field^{<|S|}[X]$ be the unique polynomial of degree less than $|S|$ with $\hat{\mathsf{Ans}}(x)=\mathsf{Ans}(x)$ for every $x\in S$ (the zero polynomial when $S$ is empty), and let $\hat{V}_S(X):=\prod_{s\in S}(X-s)$, the unique monic polynomial of degree $|S|$ vanishing on $S$.
    The \emph{quotient function} $\mathsf{Quotient}\bigl(f,S,\mathsf{Ans},\mathsf{Fill}\bigr): \evaldomain\to\field$
    is defined as follows:

    \[
    \forall x \in \evaldomain, \quad
    \mathsf{Quotient}\bigl(f, S, \mathsf{Ans}, \mathsf{Fill} \bigr)(x)
    :=
    \left\{
    \begin{array}{ll}
        \mathsf{Fill}(x)
        & \text{if } x \in S \\[6pt]
        \dfrac{f(x) - \hat{\mathsf{Ans}}(x)}{\hat{V}_S(x)}
        & \text{otherwise}
    \end{array}
    \right.
    \]

\end{definition}

Next we define the polynomial quotient operator, which quotients a polynomial relative to its output on evaluation points. The polynomial quotient is a polynomial of lower degree.

\begin{definition}\label{def:poly_quotient}
\lean{Quotienting.polyQuotient}
\leanok
    Let $\hat{f}\in\field^{<\degree}[X]$ be a polynomial and $S\subseteq\field$ be a finite set with $|S|\le\degree$, and let $\hat{V}_S(X):=\prod_{s\in S}(X-s)$, the unique monic polynomial of degree $|S|$ vanishing on $S$. The \emph{polynomial quotient} $\mathsf{PolyQuotient}(\hat{f},S)\in\field^{<d-|S|}[X]$ is defined as follows:
    \[
            \mathsf{PolyQuotient}(\hat{f},S)(X):=\frac{\hat{f}(X)-\hat{\mathsf{Ans}}(X)}{\hat{V}_S(X)}
    \]
where $\hat{\mathsf{Ans}}\in\field^{<|S|}[X]$ is the unique polynomial of degree less than $|S|$ interpolating $\hat{f}$ on $S$.
\end{definition}

The following lemma, implicit in prior works, shows that if the function is ``quotiented by the wrong value'', then its quotient is far from low-degree.

\begin{lemma}\label{lemma:quotienting}
\lean{Quotienting.quotienting}
\leanok
\uses{def:quotient,def:reed_solomon_code,def:distance_from_code,def:list_close_codewords}
    Let $f:\evaldomain\rightarrow\field$ be a function, $\degree\in\N$ be the degree parameter with $\degree\le|\evaldomain|$, $\distance\in(0,1)$
    be a distance parameter, $S\subseteq\field$ be a set with $|S|<\degree$, and $\mathsf{Ans},\mathsf{Fill}:S\rightarrow\field$ are functions. Suppose that for every $u\in\listcode(f,\degree,\distance)$ there exists $x\in S$ with $\hat{u}(x)\neq\mathsf{Ans}(x)$. Then

    \[
            \Delta(\mathsf{Quotient}(f,S,\mathsf{Ans},\mathsf{Fill}),\rscode[\field,\evaldomain,\degree-|S|])+\frac{|T|}{|\evaldomain|}>\distance,
    \]
    where $T:=\{x\in\evaldomain\cap S: \hat{\mathsf{Ans}}(x)\neq f(x)\}$.
\end{lemma}

\subsubsection{Out of domain sampling}\label{sec:out_of_domain_smpl}
\begin{lemma}\label{lemma:out_of_domain_smpl}
\lean{OutOfDomSmpl.out_of_dom_smpl_1, OutOfDomSmpl.out_of_dom_smpl_2}
\leanok
\uses{def:reed_solomon_code,def:list_decodable,def:list_close_codewords}
    Let $\field$ be a finite field and $\evaldomain\subseteq\field$ an evaluation domain with $|\evaldomain|<|\field|$, so that $\field\setminus\evaldomain$ is nonempty. Let $f:\evaldomain\rightarrow\field$ be a function, $\degree\in\N$ be a degree parameter, $s\in\N$ be a repetition parameter, $\distance\ge0$ be a distance parameter and $l\ge0$ a real list-size bound. If $\rscode[\field,\evaldomain,\degree]$ is $(\distance,l)$-list decodable, then
        \[
        \begin{array}{rcl}
        \Pr_{\substack{r_0, \ldots, r_{s-1} \gets \field \setminus \evaldomain}}
        \left[
        \begin{array}{c}
            \exists \text{ distinct } u, u' \in \listcode(f, \degree, \distance): \\
            \forall i \in [s],\ \hat{u}(r_i) = \hat{u}'(r_i)
        \end{array}
        \right]
        & \leq & \frac{l\,(l-1)}{2} \cdot \left( \frac{\degree - 1}{|\field| - |\evaldomain|} \right)^s \\
        & \leq & \left( \frac{l^2}{2} \right) \cdot \left( \frac{\degree}{|\field| - |\evaldomain|} \right)^s
        \end{array}
        \]
    Here $l-1$ is truncated at $0$, and the $r_i$ are independent and uniform on
    $\field\setminus\evaldomain$.
\end{lemma}

\subsubsection{Folding univariate functions}\label{sec:folding_uf}
STIR relies on $k$-wise folding of functions and polynomials - this is similar to prior works, although presented in a slightly different form. As shown below, folding a function preserves proximity from the Reed-Solomon code with high probability. The folding operator is based on the following fact, decomposing univariate polynomials into bivariate ones.

\begin{lemma}\label{fact:poly_folding}
\lean{Polynomial.FoldingPolynomial.foldingPolynomial,
Polynomial.FoldingPolynomial.substitution_property_of_folding_polynomial,
Polynomial.FoldingPolynomial.folding_polynomial_deg_x,
Polynomial.FoldingPolynomial.folding_polynomial_deg_y_bound,
Polynomial.FoldingPolynomial.folding_polynomial_is_unique,
Polynomial.FoldingPolynomial.folding_polynomial_deg_x_bound,
Polynomial.FoldingPolynomial.folded_poly_degree_bound}
\leanok
Given a non-constant polynomial $\hat{q}\in\field[X]$:
\begin{itemize}
        \item For every univariate polynomial $\hat{f}\in\field[X]$, there exists a unique bivariate polynomial $\hat{Q}\in\field[X,Y]$ with:
        \[
          \polydeg_X(\hat{Q}) := \left\lfloor \frac{\polydeg(\hat{f})}{\polydeg(\hat{q})} \right\rfloor,\quad \polydeg_Y(\hat{Q}) < \polydeg(\hat{q})
        \]
        such that $\hat{f}(Z)=\hat{Q}(\hat{q}(Z),Z)$. Observe that if $\polydeg(\hat{f})<t\cdot\polydeg(\hat{q})$ then
        $\polydeg_X(\hat{Q})<t$.

        \item For every $\hat{Q}[X,Y]$ with $\polydeg_X(\hat{Q})<t$ and $\polydeg_Y(\hat{Q})<\polydeg(\hat{q})$, the polynomial $\hat{f}(Z)=\hat{Q}(\hat{q}(Z),Z)$ has degree $\polydeg(\hat{f})<t\cdot\polydeg(\hat{q})$.
\end{itemize}
\end{lemma}

Moreover, $\hat{Q}$ can be computed efficiently given $\hat{f}$ and $\hat{q}$; the formalization
defines it by repeated division by $\hat{q}$.

Below, we define folding of a polynomial followed by folding of a function.
\begin{definition}\label{def:poly_folding}
\lean{Polynomial.FoldingPolynomial.polyFold,
Polynomial.FoldingPolynomial.polyFold_natDegree_le}
\leanok
\uses{fact:poly_folding}
    Given a polynomial $\hat{f}\in\field[X]$, a folding parameter $k\in\N$ with $k\ge1$ and $r\in\field$, we define a polynomial $\mathsf{PolyFold}(\hat{f},k,r)\in\field[X]$ as follows. Let $\hat{Q}[X,Y]$ be the bivariate polynomial derived from $\hat{f}$ using Fact~\ref{fact:poly_folding} with the non-constant polynomial $\hat{q}(X):=X^k$. Then $\mathsf{PolyFold}(\hat{f},k,r)(X):=\hat{Q}(X,r)$. Its degree is at most $\lfloor\polydeg(\hat{f})/k\rfloor$.
\end{definition}

\begin{definition}\label{def:fn_folding}
\lean{ProximityGap.foldWord}
\leanok
Let $f:\evaldomain\rightarrow\field$ be a function, $k\in\N$ a folding parameter and $\alpha\in\field$. For every $x\in{\evaldomain}^k$, let $\hat{p}_x\in\field^{<k}[X]$ be the polynomial where $\hat{p}_x(y)=f(y)$ for every $y\in{\evaldomain}$ such that $y^k=x$. We define $\mathsf{Fold}(f,k,\alpha):{\evaldomain}^k\rightarrow\field$ as follows.
\[
    \mathsf{Fold}(f,k,\alpha)(x):=\hat{p}_x(\alpha).
\]
In order to compute $\mathsf{Fold}(f,k,\alpha)(x)$ it suffices to interpolate the $k$ values $\{f(y):y\in\evaldomain\text{ s.t. }y^k=x\}$ into the polynomial $\hat{p}_x$ and evaluate this polynomial at $\alpha$.
The formalization takes $\evaldomain$ to be a smooth coset FFT domain of size $2^n$ and a folding
parameter of the form $2^k$.
\end{definition}

The following lemma shows that the distance of a function is preserved under folding. If a functions $f$ has distance $\distance$ to a Reed-Solomon code then, with high probability over the choice of folding randomness, its folding also has a distance of $\distance$ to the ``$k$-wise folded'' Reed-Solomon code.

\begin{lemma}\label{lemma:folding}
\lean{ProximityGap.folding_preserves_distance}
\leanok
\uses{def:reed_solomon_code,def:fn_folding,def:distance_from_code}
    Let $\evaldomain$ be a smooth coset FFT domain of size $2^n$, let $\degree=2^{d}$ and let the
    folding parameter be $2^{k}$ with $1\le k\le d\le n$. Let $\rate:=\frac{\degree}{|\evaldomain|}$.
    For every function $f:\evaldomain\rightarrow\field$ and distance parameter
    $\distance\in(0,\min\{\Delta(f,\rscode[\field,\evaldomain,\degree]),1-{\mathsf{B}}^*(\rate)\})$,
    \[
        \Pr_{r^{\mathsf{fold}}\gets\field}[\Delta(\mathsf{Fold}(f,2^k,r^{\mathsf{fold}}),\rscode[\field,{\evaldomain}^{2^k}, \degree/2^k])\le\distance]\le\err^*(\degree/2^k,\rate,\distance,2^k).
    \]
    Above, ${\mathsf{B}}^*$ and ${\err}^*$ are the proximity bound and error (respectively) described in Section~\ref{sec:proximity_gap}.
\end{lemma}

\subsubsection{Combine functions of varying degrees}\label{sec:combine}
We show a new method for combining functions of varying degrees with minimal proximity require- ments using geometric sums. We begin by recalling a fact about geometric sums.

\begin{lemma}\label{fact:geometric_sum}
\lean{Combine.geometric_sum_units}
\leanok
    Let $\field$ be a field, $r\in\field^\times$ be a nonzero field element, $a\in\N$ be a natural number. Then
    \[
    \sum_{i=0}^{a} r^i :=
    \left\{
    \begin{array}{ll}
        \displaystyle\left(\frac{1 - r^{a+1}}{1 - r}\right)
        & \text{if } r \neq 1 \\[6pt]
        a + 1
        & \text{if } r = 1
    \end{array}
    \right.
    \]

\end{lemma}

\begin{definition}\label{def:combine}
\lean{Combine.combine}
\leanok
    Given target degree $\degree^*\in\N$, shifting parameter $r\in\field$, functions $f_0,\ldots,f_{m-1}:\evaldomain\rightarrow\field$, and degrees $0\leq \degree_0,\ldots,\degree_{m-1}\leq {\degree}^*$, we define $\combine(\degree^*,r,(f_0,\degree_0),\ldots,(f_{m-1},\degree_{m-1})):\evaldomain\rightarrow\field$ as follows:
        \[
        \begin{array}{rcl}
        \combine(\degree^*, r, (f_0, \degree_0), \ldots, (f_{m-1}, \degree_{m-1}))(x)
        & := & \displaystyle\sum_{i=0}^{m-1} r_i \cdot f_i(x) \cdot \left( \sum_{l=0}^{\degree^* - \degree_i} (r \cdot x)^l \right) \\[10pt]
        & = & \left\{
        \begin{array}{ll}
            \displaystyle\sum_{i=0}^{m-1} r_i \cdot f_i(x) \cdot \left( \dfrac{1 - (x r)^{\degree^* - \degree_i + 1}}{1 - x r} \right)
            & \text{if } x \cdot r \neq 1 \\[8pt]
            \displaystyle\sum_{i=0}^{m-1} r_i \cdot f_i(x) \cdot (\degree^* - \degree_i + 1)
            & \text{if } x \cdot r = 1
        \end{array}
        \right.
        \end{array}
        \]

Above, $r_i:=r^{i+\sum_{j<i}(\degree^*-\degree_j)}$.
\end{definition}

\begin{definition}\label{def:deg_corr}
\lean{Combine.degCor}
\leanok
\uses{def:combine,def:reed_solomon_code,def:distance_from_code}
    Given target degree $\degree^*\in\N$, shifting parameter $r\in\field$, function $f:\evaldomain\rightarrow\field$, and degree $0\leq\degree\leq\degree^*$, we define $\degcorr(\degree^*,r,f,\degree)$ as follows.
    \[
    \begin{array}{rcl}
    \degcorr(\degree^*, r, f, \degree)(x)
    & := & \displaystyle f(x) \cdot \left( \sum_{l=0}^{\degree^* - \degree} (r \cdot x)^l \right) \\[10pt]
    & = & \left\{
    \begin{array}{ll}
        \displaystyle f(x) \cdot \left( \dfrac{1 - (x r)^{\degree^* - \degree + 1}}{1 - x r} \right)
        & \text{if } x \cdot r \neq 1 \\[8pt]
        \displaystyle f(x) \cdot (\degree^* - \degree + 1)
        & \text{if } x \cdot r = 1
    \end{array}
    \right.
    \end{array}
    \]


(Observe that $\degcorr(\degree^*,r,f,\degree)=\combine(\degree^*,r,(f,\degree))$.)
\end{definition}

Below it is shown that combining multiple polynomials of varying degrees can be done as long as the proximity error is bounded by $(\min{\{1-{\mathsf{B}}^*(\rate),1-\rate-1/|\evaldomain|\}})$.

\begin{lemma}\label{lemma:combine}
\lean{Combine.combine_theorem}
\leanok
\uses{def:reed_solomon_code,def:combine,def:deg_corr,thm:proximity_gap}
    Let $\degree^*\ge1$ be a target degree, $f_0,\ldots,f_{m-1}:\evaldomain\rightarrow\field$ be functions, $0\leq \degree_0,\ldots,\degree_{m-1}\leq \degree^*$ be degrees, $\distance\in(0,\min{\{1-{\mathsf{B}}^*(\rate),1-\rate-1/|\evaldomain|\}})$ be a distance parameter, where $\rate=\degree^*/|\evaldomain|$. If
    \[
        \Pr_{r\gets\field}[\Delta(\combine(\degree^*,r,(f_0,\degree_0),\ldots,(f_{m-1},\degree_{m-1})),\rscode[\field,\evaldomain,\degree^*])\le\distance]>\err^*(\degree^*,\rate,\distance,m\cdot(\degree^*+1)-\sum_{i=0}^{m-1}\degree_i),
    \]
    then there exists $S\subseteq \evaldomain$ with $|S|\geq(1-\distance)\cdot|\evaldomain|$, and
    \[
        \forall i\in\{0,\ldots,m-1\},\exists u\in\rscode[\field,\evaldomain,\degree_i], f_i(S)=u(S).
    \]
    Note that this implies $\Delta(f_i,\rscode[\field,\evaldomain,\degree_i])\le\distance$ for every $i$. Above, ${\mathsf{B}}^*$ and ${\err}^*$ are the proximity bound and error (respectively) described in the proximity gap theorem.
\end{lemma}

\subsection{STIR Section 5 targets}

Construction 5.2, Lemma 5.4 and Theorem 5.1 of \cite{ACFY24stir} remain unformalized.
The statements below record the paper's intended bounds. ArkLib formalizes the construction's
parameter and domain conditions in \texttt{StirIOP.ParamConditions}, together with the codeword
completeness relation and strict soundness relation, but has no concrete STIR prover and verifier.
A protocol must be fixed by Construction 5.2 before its distances and list sizes are chosen.

\begin{theorem}[STIR Main Theorem, unformalized target]\label{thm:stir}
\uses{def:reed_solomon_code,lemma:rnd_by_rnd_soundness}
    Consider the following ingredients:
    \begin{itemize}
        \item A security parameter $\lambda\in\N$.
        \item A Reed-Solomon code $\rscode[\field,\evaldomain,\degree]$ with $\rate:=\frac{\degree}{|\evaldomain|}$ where $\degree$ is a power of $2$, and $\evaldomain$ is a smooth domain.
        \item A proximity parameter $\distance\in(0,1-1.05\cdot\sqrt{\rate})$.
        \item A folding parameter $k\in\N$ that is power of $2$ with $k\geq 4$.
    \end{itemize}
If $|\field|=\Omega(\frac{\lambda\cdot2^\lambda\cdot\degree^2\cdot{|\evaldomain|}^{3.5}}{\log(1/\rate)})$, there is a public-coin IOPP for $\rscode[\field,\evaldomain,\degree]$ with the following parameters:
\begin{itemize}
    \item Round-by-round soundness error $2^{-\lambda}$.
    \item Round complexity: $M:=O(\log_k{\degree})$.
    \item Proof length: $|\evaldomain|+O_k(\log{\degree})$.
    \item Query complexity to the input: $\frac{\lambda}{-\log{(1-\distance)}}$.
    \item Query complexity to the proof strings: $O_k(\log{\degree}+\lambda\cdot\log{\Big(\frac{\log{\degree}}{\log{1/\rate}}\Big)})$.
\end{itemize}
\end{theorem}

\begin{remark}[Formalization target]
The exponent $3.5$ of the field size is the one of the January 2025 revision of ePrint 2024/390
\cite{ACFY24stir} (Theorem 5.1 and Appendix C.1). The constants in the $\Omega$ and $O$ bounds must
be uniform in the protocol parameters; $O_k$ constants may depend on $k$.
A formalization must establish the concrete protocol's round, proof-length and both query bounds,
counting $k$ grouped points as one symbol as in the paper. Numerical error and proof-length bounds
for an arbitrary protocol permit an exhaustive exact-membership verifier and do not characterize
this theorem. Appendix C.1 also includes a field-size dependence on
$1/(-\log(1-\distance'))$, absent from the displayed Theorem 5.1; the parameter analysis must
resolve this discrepancy.
\end{remark}

\begin{lemma}[Round-by-round soundness, unformalized target]\label{lemma:rnd_by_rnd_soundness}
\uses{def:reed_solomon_code,def:list_decodable,lemma:folding,lemma:out_of_domain_smpl,lemma:quotienting,lemma:combine}
    Consider $(\field,M,\degree,k_0,\ldots,k_M,\evaldomain_0,\ldots,\evaldomain_M,t_0,\ldots,t_M,s)$ and for every $i\in\{0,\ldots,M\}$, let $\degree_i:=\frac{\degree}{\prod_{j<i}k_j}$ and $\rate_i:=\degree_i/|\evaldomain_i|$, where, as in Construction~5.2 of \cite{ACFY24stir}, $\degree$ and $k_0,\ldots,k_M$ are powers of $2$ with $\degree\ge\prod_i k_i$, each $\evaldomain_i$ is a smooth domain with $|\evaldomain_i|>\degree_i$, and $t_i+1\le\degree_{i+1}$ for every $i<M$. For every $f\notin\rscode[\field,\evaldomain_0,\degree_0]$ and every $\distance_0,\ldots,\distance_M$ where
    \begin{itemize}
        \item $\distance_0\in(0,\Delta(f,\rscode[\field,\evaldomain_0,\degree_0])]\cap(0,1-{\mathsf{B}}^*(\rate_0))$
        \item for every $0<i\leq M$: $\distance_i\in(0,\min{\{1-\rate_i-\frac{1}{|\evaldomain_i|},1-{\mathsf{B}^*(\rate_i)}\}})$, and
        \item for every $0<i\leq M$: $\rscode[\field,\evaldomain_i,\degree_i]$ is $(\distance_i,l_i)$-list decodable,
    \end{itemize}
    The STIR IOPP of Construction~5.2 with these fixed parameters has round-by-round soundness error $(\epsilon^{\mathsf{fold}},\epsilon^{\mathsf{out}}_1,\epsilon^{\mathsf{shift}}_1,\ldots,\epsilon^{\mathsf{out}}_M,\epsilon^{\mathsf{shift}}_M,\epsilon^{\mathsf{fin}})$ where:
    \begin{itemize}
        \item $\epsilon^{\mathsf{fold}}\leq\err^*(\degree_0/k_0,\rate_0,\distance_0,k_0)$.
        \item $\epsilon^{\mathsf{out}}_i\leq\frac{l^2_i}{2}\cdot{\big(\frac{\degree_i}{|\field|-|\evaldomain_i|}\big)}^s$
        \item $\epsilon^{\mathsf{shift}}_i\leq {(1-\distance_{i-1})}^{t_{i-1}}+\err^*(\degree_i,\rate_i,\distance_i,t_{i-1}+s)+\err^*(\degree_i/k_i,\rate_i,\distance_i,k_i)$.
        \item $\epsilon^{\mathsf{fin}}\leq{(1-\distance_M)}^{t_M}$.
    \end{itemize}
    Above, ${\mathsf{B}}^*$ and ${\err}^*$ are the proximity bound and error (respectively) described in Section~\ref{sec:proximity_gap}.
\end{lemma}

\begin{remark}[Formalization target]
The protocol must be determined by Construction~5.2 before the distances $\distance_i$ and list
sizes $l_i$ are chosen. Its completeness relation consists of codewords, while its soundness
language is the strict distance relation: the bad inputs include distance exactly $\distance_0$.
Each component of the displayed error vector must bound its corresponding challenge, including
final verification when $M=0$. Construction~5.2 uses $2M+1$ interactive rounds and a randomized
final decision; a faithful representation must retain the prescribed uniform sampling from
$\field\setminus\evaldomain_i$ and folded evaluation domains. The current fixed-length field-vector
challenge representation does not in general realize these distributions.
\end{remark}
